\chapter{UI Service: User Interface State Management}
\label{chap:ui-service}

This chapter introduces the \texttt{UiService}, the central component responsible for managing the application's user interface state. Understanding this service is fundamental, as it dictates what the user sees and how the application responds to their interactions.

\section{Problem Statement}
\label{sec:ui-problem}

Consider a common user action: \textbf{switching tabs}. The application has different tabs such as ``Build'', ``Simulation'', and ``Analyze''. When the user clicks on the ``Simulation'' tab:

\begin{enumerate}
    \item The clicked button must inform the rest of the application: ``We are now in Simulation mode.''
    \item The main drawing area must update to show simulation-specific elements (e.g., animated token flow) instead of editing tools.
    \item Other UI components (sidebars, status messages) must reflect the new ``Simulation'' context.
\end{enumerate}

Without a centralized state management approach, every component would need to independently track the current tab, leading to:
\begin{itemize}
    \item \textbf{Inconsistency}: Components might display conflicting states
    \item \textbf{Tight coupling}: Components must directly communicate with each other
    \item \textbf{Maintenance burden}: Changes require updates across multiple files
\end{itemize}

The \texttt{UiService} solves this by being the \textbf{single source of truth} for what the user interface is currently displaying or doing.

\section{Conceptual Overview}
\label{sec:ui-concept}

Think of the \texttt{UiService} as the application's \textbf{control panel} or \textbf{dashboard}. It maintains critical information about the UI state:

\begin{itemize}
    \item \textbf{Active tab}: Which workflow is currently visible (Build, Simulation, Analyze, etc.)
    \item \textbf{Selected tool}: Which editing tool is active (Add Place, Move, Delete, etc.)
    \item \textbf{Simulation speed}: The playback rate for animation
    \item \textbf{Editor format}: Whether the Code tab shows JSON or PNML
\end{itemize}

Other components \textit{subscribe} to the \texttt{UiService}---similar to subscribing to a notification channel. When the service announces a change (e.g., ``The active tab is now Simulation!''), all subscribed components automatically react, ensuring the application remains responsive and consistent.

\section{State Definitions}
\label{sec:ui-states}

The application defines its UI states using TypeScript enums, providing clear, human-readable names instead of magic numbers.

\begin{lstlisting}[language=TypeScript, caption={UI state enums (src/app/tr-enums/ui-state.ts)}]
/**
 * Represents the currently active tab in the application.
 * Each tab corresponds to a different workflow/mode.
 */
export enum TabState {
    Build,       // Graphical editing of Petri net elements
    Simulation,  // Playback and animation of firing sequences
    Save,        // Export options (PNG, SVG, JSON, PNML)
    Code,        // Text-based editing of PNML/JSON
    Analyze,     // Place invariants and analysis results
    Offshore,    // Domain-specific parameter configuration
}

/**
 * Represents the currently selected editing tool in Build mode.
 */
export enum ButtonState {
    Blitz,       // Quick-add mode
    Place,       // Add a new place (circle)
    Transition,  // Add a new transition (rectangle)
    Move,        // Drag elements to reposition
    Arc,         // Connect places and transitions
    Delete,      // Remove elements
    // ... additional tools ...
}

/**
 * Format for the Code tab editor.
 */
export enum CodeEditorFormat {
    JSON,        // Custom JSON format
    PNML,        // Standard PNML XML format
}
\end{lstlisting}

Using enums like \texttt{TabState.Simulation} is far more readable and maintainable than remembering that ``tab state 1'' means simulation mode.

\section{The UiService Implementation}
\label{sec:ui-implementation}

The \texttt{UiService} class holds the actual state variables and provides reactive streams (Observables) that notify listeners of changes.

\begin{lstlisting}[language=TypeScript, caption={UiService core implementation (src/app/tr-services/ui.service.ts)}]
import { Injectable } from '@angular/core';
import { BehaviorSubject, Observable } from 'rxjs';
import { ButtonState, CodeEditorFormat, TabState } from '../tr-enums/ui-state';

@Injectable({
    providedIn: 'root', // Singleton: one instance for the entire application
})
export class UiService {
    // =========================================================================
    // TAB STATE MANAGEMENT
    // =========================================================================
    
    /** Private variable holding the current tab */
    private _tab: TabState = TabState.Build;
    
    /** 
     * BehaviorSubject: a special Observable that remembers its last value.
     * New subscribers immediately receive the current state.
     */
    private _tab$ = new BehaviorSubject<TabState>(this._tab);
    
    /** Public Observable for components to subscribe to */
    public tab$: Observable<TabState> = this._tab$.asObservable();

    /** Getter: read the current tab synchronously */
    public get tab(): TabState {
        return this._tab;
    }

    /** 
     * Setter: update the current tab.
     * Automatically notifies all subscribers via the BehaviorSubject.
     */
    public set tab(value: TabState) {
        if (this._tab !== value) {
            console.log(`[UiService] Tab changed: ${TabState[this._tab]} -> ${TabState[value]}`);
            this._tab = value;
            this._tab$.next(value); // Broadcast to all subscribers
        }
    }

    // =========================================================================
    // SIMULATION SPEED MANAGEMENT
    // =========================================================================
    
    private _simulationSpeed$ = new BehaviorSubject<number>(1); // Default: 1x speed
    public simulationSpeed$: Observable<number> = this._simulationSpeed$.asObservable();

    /** Update the simulation playback speed multiplier */
    public setSimulationSpeed(speed: number): void {
        this._simulationSpeed$.next(speed);
    }

    /** Get the current speed multiplier synchronously */
    public getAnimationSpeedMultiplier(): number {
        return this._simulationSpeed$.getValue();
    }

    // ... additional state management for tools, editor format, etc. ...
}
\end{lstlisting}

\subsection{Key Concepts Explained}

\begin{description}
    \item[@Injectable(\{providedIn: 'root'\})] Makes \texttt{UiService} a \textbf{singleton}---only one instance exists throughout the application. This ensures all components share the same state.
    
    \item[BehaviorSubject<T>] A special RxJS type that:
    \begin{itemize}
        \item Always holds a current value
        \item Immediately emits the current value to new subscribers
        \item Broadcasts new values to all active subscribers when \texttt{.next()} is called
    \end{itemize}
    
    \item[Getter/Setter pattern] The \texttt{get tab()} and \texttt{set tab()} methods provide a clean API. The setter automatically calls \texttt{\_tab\$.next(value)} to notify subscribers.
\end{description}

\section{Usage Patterns}
\label{sec:ui-usage}

\subsection{Changing the UI State}

When a user clicks a tab button, the \texttt{ButtonBarComponent} updates the \texttt{UiService}:

\begin{lstlisting}[language=TypeScript, caption={Setting the tab state (ButtonBarComponent)}]
import { Component } from '@angular/core';
import { UiService } from 'src/app/tr-services/ui.service';
import { TabState } from 'src/app/tr-enums/ui-state';

@Component({ /* ... */ })
export class ButtonBarComponent {
    constructor(protected uiService: UiService) {}

    /**
     * Called when user clicks a tab button.
     * Updates the UiService, which broadcasts to all subscribers.
     */
    tabClicked(tabName: string): void {
        switch (tabName.toLowerCase()) {
            case 'build':
                this.uiService.tab = TabState.Build;
                break;
            case 'simulation':
                this.uiService.tab = TabState.Simulation;
                break;
            case 'analyze':
                this.uiService.tab = TabState.Analyze;
                break;
            // ... other tabs ...
        }
    }
}
\end{lstlisting}

\subsection{Listening to UI State Changes}

Components that need to react to state changes \textit{subscribe} to the observable:

\begin{lstlisting}[language=TypeScript, caption={Subscribing to tab changes (AppComponent)}]
import { Component, OnInit, OnDestroy } from '@angular/core';
import { Subscription } from 'rxjs';
import { UiService } from './tr-services/ui.service';
import { TabState } from './tr-enums/ui-state';

@Component({ /* ... */ })
export class AppComponent implements OnInit, OnDestroy {
    private tabSubscription: Subscription;

    constructor(protected uiService: UiService) {}

    ngOnInit(): void {
        // Subscribe to tab changes
        this.tabSubscription = this.uiService.tab$.subscribe(currentTab => {
            console.log('Current tab:', TabState[currentTab]);
            
            // React to the change
            if (currentTab === TabState.Simulation) {
                this.showSimulationControls();
            } else {
                this.hideSimulationControls();
            }
        });
    }

    ngOnDestroy(): void {
        // Always unsubscribe to prevent memory leaks
        this.tabSubscription?.unsubscribe();
    }
    
    // ...
}
\end{lstlisting}

The \texttt{\$} suffix on \texttt{tab\$} is a common Angular/RxJS convention indicating that the property is an Observable (a ``stream'' of values over time).

\section{Data Flow Visualization}
\label{sec:ui-flow}

The following sequence illustrates what happens when a user switches tabs:

\begin{figure}[htb]
    \centering
    \begin{tikzpicture}[
        node distance=1.5cm,
        every node/.style={font=\small},
        actor/.style={rectangle, draw, minimum width=2cm, minimum height=0.8cm},
        service/.style={rectangle, draw, fill=blue!10, minimum width=2.5cm, minimum height=0.8cm},
        component/.style={rectangle, draw, fill=green!10, minimum width=2.5cm, minimum height=0.8cm},
    ]
        % Nodes
        \node[actor] (user) {User};
        \node[component, right=of user] (buttonbar) {ButtonBar};
        \node[service, right=of buttonbar] (uiservice) {UiService};
        \node[component, below=of uiservice] (app) {AppComponent};
        \node[component, below=of app] (petrinet) {PetriNetComponent};
        
        % Arrows
        \draw[->, thick] (user) -- node[above] {\scriptsize clicks tab} (buttonbar);
        \draw[->, thick] (buttonbar) -- node[above] {\scriptsize set tab} (uiservice);
        \draw[->, thick] (uiservice) -- node[right] {\scriptsize tab\$ emits} (app);
        \draw[->, thick] (uiservice.south) -- ++(0,-0.5) -| node[near start, right] {\scriptsize tab\$ emits} (petrinet.north);
        
    \end{tikzpicture}
    \caption{Tab switching data flow through UiService}
    \label{fig:ui-service-flow}
\end{figure}

\begin{enumerate}
    \item \textbf{User Action}: User clicks the ``Simulation'' tab button
    \item \textbf{State Update}: \texttt{ButtonBarComponent} sets \texttt{uiService.tab = TabState.Simulation}
    \item \textbf{Internal Update}: \texttt{UiService} updates \texttt{\_tab} and calls \texttt{\_tab\$.next(value)}
    \item \textbf{Broadcast}: All subscribers to \texttt{tab\$} receive the new value
    \item \textbf{Component Reactions}: 
    \begin{itemize}
        \item \texttt{AppComponent} shows the simulation timeline
        \item \texttt{PetriNetComponent} enables transition highlighting and hides build tools
    \end{itemize}
\end{enumerate}

\section{Practical Example: Tab-Aware Behavior}
\label{sec:ui-example}

The \texttt{PetriNetComponent} demonstrates how components use \texttt{tab\$} to conditionally enable features:

\begin{lstlisting}[language=TypeScript, caption={Tab-aware component behavior}]
import { Component, OnInit } from '@angular/core';
import { UiService } from 'src/app/tr-services/ui.service';
import { TabState } from 'src/app/tr-enums/ui-state';

@Component({ /* ... */ })
export class PetriNetComponent implements OnInit {
    private previousTab: TabState;

    constructor(private uiService: UiService) {}

    ngOnInit(): void {
        this.uiService.tab$.subscribe(currentTab => {
            // Leaving Simulation tab: clean up highlights and reset tokens
            if (this.previousTab === TabState.Simulation && currentTab !== TabState.Simulation) {
                this.clearAllHighlights();
                this.resetTokensToBaseline();
            }
            
            // Entering Simulation tab: save baseline marking
            if (currentTab === TabState.Simulation && this.previousTab !== TabState.Simulation) {
                this.saveBaselineMarking();
            }
            
            // Entering Build tab: enable editing tools
            if (currentTab === TabState.Build) {
                this.enableDragAndDrop();
                this.showEditingHandles();
            } else {
                this.disableDragAndDrop();
                this.hideEditingHandles();
            }
            
            this.previousTab = currentTab;
        });
    }
    
    // ... helper methods ...
}
\end{lstlisting}

This pattern ensures that:
\begin{itemize}
    \item Simulation highlights are cleared when leaving the Simulation tab
    \item Token counts are reset to their initial values
    \item Editing tools are only active in Build mode
\end{itemize}

\section{Summary}
\label{sec:ui-summary}

The \texttt{UiService} is the application's \textbf{control panel}, centralizing all UI state management:

\begin{itemize}
    \item Uses \textbf{TypeScript enums} for type-safe state definitions
    \item Employs \textbf{BehaviorSubject} for reactive state broadcasting
    \item Provides \textbf{getter/setter} methods for clean state access
    \item Enables \textbf{decoupled components} that react to state changes without direct communication
\end{itemize}

Understanding the \texttt{UiService} is essential because it controls what the user experiences. In the next chapter, we will examine the \textbf{Petri Net Elements}---the actual data objects (places, transitions, arcs) that the UI displays and manipulates.
